
\documentclass{article}
\usepackage{colortbl}
\usepackage{makecell}
\usepackage{multirow}
\usepackage{supertabular}

\begin{document}

\newcounter{utterance}

\centering \large Interaction Transcript for game `ifeval', experiment `ifeval', episode 303 with Qwen3.5{-}27B.
\vspace{24pt}

{ \footnotesize  \setcounter{utterance}{1}
\setlength{\tabcolsep}{0pt}
\begin{supertabular}{c@{$\;$}|p{.15\linewidth}@{}p{.15\linewidth}p{.15\linewidth}p{.15\linewidth}p{.15\linewidth}p{.15\linewidth}}
    \# & \multicolumn{2}{c}{Player} && \multicolumn{2}{c}{Game Master} \\
    \hline

    \theutterance \stepcounter{utterance}  
    & & & \multicolumn{4}{p{0.6\linewidth}}{
        \cellcolor[rgb]{0.9,0.9,0.9}{
            \makecell[{{p{\linewidth}}}]{
                \texttt{\tiny{[P1$\langle$GM]}}
                \texttt{Write a paragraph that lists the average length of various animal specimens from smallest to largest. Your response should contain less than 17 sentences.} \\
            }
        }
    }
     \\ \\

    \theutterance \stepcounter{utterance}  
    & \multicolumn{4}{p{0.6\linewidth}}{
        \cellcolor[rgb]{0.561,0.737,0.561}{
            \makecell[{{p{\linewidth}}}]{
                \texttt{\tiny{[P1$\rangle$GM]}}
                \texttt{The journey through the animal kingdom's size spectrum begins with the microscopic parasitic wasp *Dicopomorpha echmepterygis*, which measures a mere 0.139 millimeters in length. Slightly larger is the bumblebee bat, also known as the Kitti's hog{-}nosed bat, averaging just 29 to 33 millimeters from nose to tail. Moving up the scale, the pygmy marmoset stands as the smallest monkey, typically reaching about 14 centimeters in body length. The common house mouse follows, with an average body length of roughly 7.5 to 10 centimeters, excluding its long tail. A domestic cat usually measures between 46 and 51 centimeters from head to tail base, making it significantly larger than the rodent. The red fox, a common wild canine, averages around 45 to 90 centimeters in body length, not including its bushy tail. A gray wolf is much more imposing, with adults ranging from 105 to 160 centimeters in body length. The African elephant represents a massive leap in size, with adults standing 2.2 to 3.5 meters tall at the shoulder and reaching lengths of up to 7 meters. Even larger is the blue whale, the largest animal known to have ever lived, which can grow to an astonishing 24 to 30 meters in length. This progression from the nearly invisible wasp to the ocean{-}dwelling giant illustrates the incredible diversity of biological forms on Earth. Each species has evolved its specific size to suit its unique ecological niche and survival strategies. From the tiny insect navigating flower petals to the colossal whale migrating across oceans, size plays a critical role in their existence. The differences in scale are so vast that they often defy human intuition regarding mass and volume. Understanding these measurements helps scientists classify and study the relationships between different organisms. Ultimately, the range of animal sizes is a testament to the adaptability of life in our world.} \\
            }
        }
    }
    & & \\ \\

    \theutterance \stepcounter{utterance}  
    & & & \multicolumn{2}{p{0.3\linewidth}}{
        \cellcolor[rgb]{0.9,0.9,0.9}{
            \makecell[{{p{\linewidth}}}]{
                \texttt{\tiny{[GM$|$GM]}}
                {'length_constraints:number_sentences': {'relation': 'less than', 'num_sentences': 17}}
            }
        }
    }
    & & \\ \\

    \theutterance \stepcounter{utterance}  
    & & & \multicolumn{2}{p{0.3\linewidth}}{
        \cellcolor[rgb]{0.9,0.9,0.9}{
            \makecell[{{p{\linewidth}}}]{
                \texttt{\tiny{[GM$|$GM]}}
                \texttt{game\_result = WIN} \\
            }
        }
    }
    & & \\ \\

\end{supertabular}
}

\end{document}
